\documentclass[11pt]{article}
\usepackage[utf8]{inputenc}
\usepackage{amsmath}
\usepackage{amssymb}
\usepackage{graphicx}
\usepackage{hyperref}
\usepackage{listings}
\usepackage{xcolor}
\usepackage{geometry}
\geometry{margin=1in}

\title{Development Summary: Gaussian Process-Enhanced SINDy for Lotka-Volterra System Identification}
\author{Development Session Documentation}
\date{\today}

\begin{document}

\maketitle

\section{Introduction}

This document summarizes the development work on implementing and improving a Gaussian Process (GP)-enhanced Sparse Identification of Nonlinear Dynamics (SINDy) approach for identifying Lotka-Volterra equations from sparse, noisy data. The work involved comparing two distinct pathways: (1) direct ESINDy on sparse data, and (2) GP-enhanced ESINDy with data augmentation and analytical derivative computation.

\section{Problem Statement}

The goal was to compare two approaches for modeling Lotka-Volterra equations from sparse and noisy observations:

\textbf{Path 1:} Direct ESINDy on 6 sparse noisy data points using discrete-time formulation ($X(t+1) = f(X(t))$).

\textbf{Path 2:} GP-enhanced ESINDy that:
\begin{itemize}
    \item Fits a GP to sparse noisy data
    \item Samples additional plausible data points from the GP posterior
    \item Computes analytical derivatives from the GP
    \item Combines original and GP-sampled data
    \item Runs continuous-time ESINDy ($dX/dt = f(X)$) with an expanded library
\end{itemize}

\section{Initial Challenges}

\subsection{GP Fitting Issues}

Initially, the GP mean predictions were producing straight lines (R² ≈ 0.0000), indicating a complete failure to capture the underlying dynamics. This was attributed to:
\begin{itemize}
    \item Extremely sparse data (only 6 points)
    \item High noise levels (10\% Gaussian noise)
    \item Poor hyperparameter initialization
    \item Inappropriate kernel or basis function choices
\end{itemize}

\subsection{Data Generation Improvements}

To improve GP fitting, we increased the data density:
\begin{itemize}
    \item \textbf{Original approach:} 6 unique time points
    \item \textbf{Improved approach:} 18 total data points (3 replicates at each of 6 unique time points)
    \item This provided better statistical power for GP hyperparameter estimation
    \item Path 1 was modified to use averaged data (6 points) for consistency
\end{itemize}

\section{Kernel Selection and Evolution}

\subsection{Rational Quadratic Kernel (Initial Attempt)}

Initially, we attempted to use the Rational Quadratic (RQ) kernel, which is a scale mixture of squared exponential kernels and can model functions with varying smoothness:

\begin{equation}
k(x, x') = \sigma^2 \left(1 + \frac{(x-x')^2}{2\alpha l^2}\right)^{-\alpha}
\end{equation}

\textbf{Issues encountered:}
\begin{itemize}
    \item GP optimization converged to a poor local minimum with length scale $l \approx 0.0000$
    \item This caused numerical instability in analytical derivative computation
    \item The derivative formula involves division by $l^2$, leading to near-infinite values
    \item Analytical derivatives produced essentially zero values (std = 0.0002)
    \item This resulted in all-zero coefficients in SINDy regression
\end{itemize}

\textbf{Why this happened:} With sparse, noisy data, the GP optimization found a degenerate solution where the length scale collapsed to near-zero, essentially treating each data point as independent.

\subsection{Matern32 Kernel (Intermediate Attempt)}

We briefly tried the Matern32 kernel as an alternative:

\begin{equation}
k(r) = \sigma^2\left(1 + \frac{\sqrt{3}r}{l}\right)\exp\left(-\frac{\sqrt{3}r}{l}\right), \quad r = |x - x'|
\end{equation}

This kernel is less smooth than squared exponential but more flexible than RQ in some contexts. However, we did not fully evaluate this option before moving to squared exponential.

\subsection{Squared Exponential Kernel (Final Choice)}

Following the approach in Wang \& Barber (2014) \cite{wang2014}, we adopted the squared exponential (SE) kernel:

\begin{equation}
k(x, x') = \sigma^2 \exp\left(-\frac{(x-x')^2}{2l^2}\right)
\end{equation}

\textbf{Rationale:}
\begin{itemize}
    \item Well-established for smooth functions (appropriate for ODE trajectories)
    \item Stable analytical derivatives
    \item Proven effective in GP-ODE gradient matching literature
    \item Simpler parameter space (2 parameters: $l$, $\sigma_f$) compared to RQ (3 parameters)
    \item Less prone to degenerate solutions with sparse data
\end{itemize}

\section{Basis Function Selection}

\subsection{Linear Basis Function (Initial Attempt)}

Initially, we used a linear basis function, which allows the GP to model trends:

\begin{equation}
\mu(x) = h(x)^T \beta, \quad h(x) = [1, x]^T
\end{equation}

\textbf{Issues:}
\begin{itemize}
    \item For variable Y (predator), GP fit quality was extremely poor (R² = 0.0221)
    \item The GP mean became nearly constant, producing zero derivatives
    \item This was likely due to overfitting to the linear trend rather than the kernel structure
\end{itemize}

\subsection{Constant Basis Function (Final Choice)}

We switched to a constant basis function:

\begin{equation}
\mu(x) = \beta_0
\end{equation}

\textbf{Rationale:}
\begin{itemize}
    \item Matches the approach in Wang \& Barber (2014)
    \item Forces the GP kernel to capture all structure (no linear trend assumption)
    \item Improved Y GP fit quality significantly
    \item Both X and Y derivatives became non-zero and reasonable
    \item More appropriate for oscillatory dynamics where linear trends are not expected
\end{itemize}

\section{Analytical GP Derivative Computation}

\subsection{Mathematical Foundation}

For a GP with mean function $\mu(x)$ and covariance function $k(x, x')$, the derivative is itself a GP with:
\begin{itemize}
    \item Mean: $\mu'(x) = \frac{\partial \mu(x)}{\partial x}$
    \item Covariance: $k'(x, x') = \frac{\partial^2 k(x, x')}{\partial x \partial x'}$
\end{itemize}

For prediction at new points $X_*$, the derivative mean is:
\begin{equation}
\mu'(X_*) = K'(X_*, X) K(X, X)^{-1} y
\end{equation}

where $K'(X_*, X)$ is the cross-covariance matrix of derivatives.

\subsection{Implementation for Squared Exponential Kernel}

For the SE kernel, the derivative formulas are:
\begin{align}
\frac{\partial k(x, x')}{\partial x} &= -\frac{\sigma^2}{l^2}(x - x') \exp\left(-\frac{(x-x')^2}{2l^2}\right) \\
\frac{\partial^2 k(x, x')}{\partial x \partial x'} &= \frac{\sigma^2}{l^2}\left(1 - \frac{(x-x')^2}{l^2}\right) \exp\left(-\frac{(x-x')^2}{2l^2}\right)
\end{align}

\subsection{Finite Differences Fallback}

To ensure robustness, we implemented a finite differences fallback mechanism:

\begin{lstlisting}[language=Matlab, basicstyle=\small]
if std(deriv_gp) < 0.01 || isnan(std(deriv_gp))
    % Fallback to finite differences on GP mean
    dt = mean(diff(X_new));
    deriv_gp = gradient(gp_mean, dt);
end
\end{lstlisting}

\textbf{Rationale:}
\begin{itemize}
    \item Detects when analytical derivatives are unreliable (near-zero variance)
    \item Provides a stable alternative when kernel parameters are problematic
    \item Finite differences on the GP mean are still better than finite differences on raw noisy data
    \item Ensures the pipeline always produces usable derivatives
\end{itemize}

\section{GP Sampling Strategy}

\subsection{Posterior Sampling}

We sample full curves from the GP posterior distribution rather than just adding noise to the mean:

\begin{equation}
f_* \sim \mathcal{N}(\mu_*, \Sigma_*)
\end{equation}

where $\mu_*$ and $\Sigma_*$ are the GP predictive mean and covariance at new points.

\textbf{Benefits:}
\begin{itemize}
    \item Respects the full covariance structure of the GP
    \item Provides realistic variability that accounts for uncertainty
    \item Each sampled curve is a valid GP realization
    \item More principled than adding independent noise to the mean
\end{itemize}

\subsection{Data Combination Strategy}

We combine GP mean and samples using an interleaving strategy:
\begin{itemize}
    \item Even indices: Use GP mean (smooth, denoised estimate)
    \item Odd indices: Use GP samples (captures variability)
    \item Original data points: Preserved at their exact time locations
\end{itemize}

This provides both the smooth mean estimate and the variability from sampling, giving ESINDy a richer dataset while maintaining connection to the original observations.

\section{Code Structure and Implementation}

\subsection{Key Functions}

\textbf{\texttt{gp\_derivative\_analytical.m:}}
\begin{itemize}
    \item Computes analytical derivatives for multiple kernel types
    \item Supports Squared Exponential, Matern32, and Rational Quadratic
    \item Includes fallback to finite differences
    \item Handles kernel-specific parameter extraction
\end{itemize}

\textbf{\texttt{sample\_gp\_posterior.m:}}
\begin{itemize}
    \item Samples full curves from GP posterior
    \item Supports multiple kernel types
    \item Computes full covariance matrices for proper sampling
\end{itemize}

\textbf{\texttt{lotka\_volterra\_comparison.m:}}
\begin{itemize}
    \item Main comparison script
    \item Implements both Path 1 and Path 2
    \item Handles data generation, GP fitting, sampling, and ESINDy execution
    \item Includes comprehensive diagnostics and visualization
\end{itemize}

\section{Results and Current Status}

\subsection{GP Fitting Quality}

With squared exponential kernel and constant basis:
\begin{itemize}
    \item X (Prey) GP: R² = 0.9836 (excellent fit)
    \item Y (Predator) GP: R² improved significantly (from 0.0221 with linear basis)
    \item Both derivatives are non-zero and reasonable
    \item Kernel parameters are stable (no near-zero length scales)
\end{itemize}

\subsection{Derivative Quality}

\begin{itemize}
    \item X derivative: std = 4.30 (non-zero, captures dynamics)
    \item Y derivative: std = 2.24 (non-zero, captures dynamics)
    \item Both derivatives computed analytically from GP
    \item Finite differences fallback available but not needed with current setup
\end{itemize}

\subsection{Model Performance}

Path 1 (Direct ESINDy):
\begin{itemize}
    \item Uses 6 averaged data points
    \item Discrete-time formulation
    \item Library: [1, x, y, x², y², xy] (6 functions)
    \item RMSE: X = 0.16, Y = 0.11 (Average = 0.14)
\end{itemize}

Path 2 (GP-enhanced ESINDy):
\begin{itemize}
    \item Uses 50 combined points (6 original + 44 GP-sampled)
    \item Continuous-time formulation with GP derivatives
    \item Expanded library: [1, x, y, x², y², xy, x³, y³, x²y, xy²] (10 functions)
    \item Currently experiencing model instability (high RMSE)
    \item Likely due to overfitting with expanded library or integration issues
\end{itemize}

\section{Lessons Learned}

\subsection{What Worked}

\begin{enumerate}
    \item \textbf{Squared Exponential Kernel:} Stable, well-behaved, appropriate for smooth ODE trajectories
    \item \textbf{Constant Basis Function:} Better than linear for oscillatory dynamics
    \item \textbf{Analytical Derivatives:} More principled than finite differences, especially with sparse data
    \item \textbf{Posterior Sampling:} Provides realistic data augmentation
    \item \textbf{Finite Differences Fallback:} Essential safety mechanism for robustness
    \item \textbf{Increased Data Density:} 18 points (3 replicates) much better than 6 points
\end{enumerate}

\subsection{What Didn't Work}

\begin{enumerate}
    \item \textbf{Rational Quadratic Kernel:} Prone to degenerate solutions (length scale → 0) with sparse data
    \item \textbf{Linear Basis Function:} Overfitted to trends, poor fit for oscillatory Y variable
    \item \textbf{Very Sparse Data:} 6 points insufficient for reliable GP hyperparameter estimation
    \item \textbf{Expanded Library:} May cause overfitting; needs regularization or model selection
\end{enumerate}

\subsection{Key Insights}

\begin{enumerate}
    \item Kernel choice is critical for sparse data—simpler kernels (SE) more robust than complex ones (RQ)
    \item Basis function choice should match problem structure (constant better than linear for oscillations)
    \item GP hyperparameter optimization can find degenerate solutions—need diagnostics and fallbacks
    \item Analytical derivatives are superior but require stable kernel parameters
    \item More data (even with replicates) significantly improves GP fitting
    \item Model complexity (library size) must be balanced with data availability
\end{enumerate}

\section{Future Directions}

\subsection{Potential Improvements}

\begin{enumerate}
    \item \textbf{Hyperparameter Constraints:} Set minimum length scale bounds to prevent degenerate solutions
    \item \textbf{Model Selection:} Implement cross-validation or information criteria for library selection
    \item \textbf{Regularization:} Add stronger sparsity penalties for expanded library
    \item \textbf{Uncertainty Quantification:} Propagate GP uncertainty through to SINDy coefficients
    \item \textbf{Alternative Kernels:} Explore periodic kernels for oscillatory dynamics
    \item \textbf{Multi-scale Approaches:} Use composite kernels (e.g., SE + periodic) for complex dynamics
\end{enumerate}

\section{Conclusion}

This development session successfully implemented a GP-enhanced SINDy pipeline for Lotka-Volterra system identification. Through iterative refinement, we identified that:

\begin{itemize}
    \item Squared exponential kernel with constant basis function provides stable, high-quality GP fits
    \item Analytical GP derivatives are computable and reliable with appropriate kernel choices
    \item Data augmentation through GP posterior sampling provides principled expansion of sparse datasets
    \item Robust fallback mechanisms are essential for production code
\end{itemize}

The current implementation provides a solid foundation for comparing direct ESINDy with GP-enhanced approaches, with clear paths for further improvement in model selection and regularization.

\section{References}

\begin{itemize}
    \item Wang, Y., \& Barber, D. (2014). Gaussian Processes for Bayesian Estimation in Ordinary Differential Equations. \textit{ICML}.
    \item Brunton, S. L., Proctor, J. L., \& Kutz, J. N. (2016). Discovering governing equations from data by sparse identification of nonlinear dynamical systems. \textit{PNAS}, 113(15), 3932-3937.
    \item Rasmussen, C. E., \& Williams, C. K. I. (2006). \textit{Gaussian Processes for Machine Learning}. MIT Press.
\end{itemize}

\end{document}

